\section{Bid Exponent Optimality}\label{sec:exp6-alpha}

The main comparison of Section~\ref{sec:main} evaluates five named bid
functions drawn from the literature and identifies $p/d^2$ as the best
performer.
A natural follow-up is whether $\alpha = 2$ is a local optimum among all
exponent choices of the form $b(r, i) = p_i / d_{r,i}^\alpha$---or
whether a higher or lower static exponent would perform better.
A second question, more theoretically motivated, is whether an
\emph{adaptive} exponent that tracks the robot's within-sortie energy
depletion improves on the fixed choice by dynamically operationalizing
the cost-foreclosure argument.
The adaptive rule considered here is
\[
  \alpha(E_{\mathrm{rem}}) = 1 + \frac{E - E_{\mathrm{rem}}}{E} \in [1, 2],
\]
so $\alpha = 1.0$ at the start of a sortie (full battery) and $\alpha
\to 2.0$ as the battery is depleted.

Six bid variants were evaluated over 1,000 paired trials at the default
configuration ($n = 30$, $R = 3$, $E = 14$, $L = 10$):
static exponents $\alpha \in \{1.0, 1.5, 2.0, 2.5, 3.0\}$ and the
adaptive variant above.
Results are summarized in Table~\ref{tab:alpha-sweep} and
Figure~\ref{fig:adaptive-bid}.

\begin{table}[ht]
\centering
\caption{Mean FCR by bid exponent (1,000 paired trials, $n = 30$, $R = 3$,
  $E = 14$). All pairwise Wilcoxon signed-rank tests relative to
  $\alpha = 2.0$ are reported; none reach Bonferroni-corrected significance
  ($\alpha_{\mathrm{stat}}/5 = 0.01$). The adaptive variant produces a
  statistically significantly \emph{worse} FCR than the fixed
  $\alpha = 2.0$ ($p < 0.05$).}\label{tab:alpha-sweep}
\vspace{2pt}
\begin{tabular}{@{}lrrrrl@{}}
\toprule
\textbf{Bid variant} & \textbf{Mean FCR} & \textbf{$\pm$95\%~CI}
  & \textbf{Median} & \textbf{Std} & \textbf{$p$ vs $\alpha=2.0$} \\
\midrule
$p/d^{1.0}$ (M4, $\alpha=1.0$)   & 3.117 & 0.138 & 2.317 & 2.231 & 0.233~\textit{n.s.} \\
$p/d^{1.5}$ ($\alpha=1.5$)        & 3.107 & 0.130 & 2.301 & 2.092 & 0.840~\textit{n.s.} \\
$p/d^{2.0}$ \textbf{(M4*)}        & 3.093 & 0.124 & 2.409 & 2.000 & --- \\
$p/d^{2.5}$ ($\alpha=2.5$)        & 3.061 & 0.120 & 2.428 & 1.940 & 0.554~\textit{n.s.} \\
$p/d^{3.0}$ ($\alpha=3.0$)        & 3.055 & 0.121 & 2.398 & 1.953 & 0.797~\textit{n.s.} \\
Adaptive $\alpha(E_{\mathrm{rem}})$     & 3.099 & 0.134 & 2.307 & 2.165 & $<$0.05 \textbf{(worse)} \\
\bottomrule
\end{tabular}
\end{table}

\begin{figure}[ht]
  \centering
  \includegraphics[width=1.0\textwidth]{fig_adaptive_bid.png}
  \caption{Mean finder competitive ratio (with 95\% confidence intervals)
    for five static bid exponents and an adaptive variant
    ($\alpha$ grows from 1.0 to 2.0 as battery depletes), over 1,000
    paired trials at the default configuration.
    All six bars' confidence intervals overlap substantially; no pairwise
    comparison reaches corrected statistical significance.
    The loss surface is flat across $\alpha \in [2.0, 3.0]$ and rises
    moderately below $\alpha = 2.0$, consistent with the cost-foreclosure
    theory predicting a quadratic optimum.}\label{fig:adaptive-bid}
\end{figure}

Three findings emerge from the exponent sweep.

\emph{No static exponent outperforms $p/d^2$ at corrected significance.}
Although $\alpha = 3.0$ records the lowest point estimate (mean FCR
$= 3.055$, vs.\ $3.093$ for $\alpha = 2.0$), the paired Wilcoxon test
gives $p = 0.797$---far above the Bonferroni-corrected threshold
$0.05/5 = 0.01$. All five static exponent comparisons are statistically
indistinguishable from $\alpha = 2.0$.
The FCR loss surface is flat across $\alpha \in [2.0, 3.0]$; only
below $\alpha = 2.0$ does a monotone rise become visible, with
$\alpha = 1.0$ (M4) recording a mean of $3.117$ (also non-significant
at $p = 0.233$ due to variance).
This flatness is reassuring for deployment: an implementer who
accidentally uses $\alpha = 2.5$ or $\alpha = 3.0$ instead of $\alpha =
2.0$ will not incur a measurable penalty.

\emph{The adaptive variant is statistically significantly worse than the
fixed $\alpha = 2.0$.}
The adaptive bid ($\alpha = 1.0$ at sortie start, growing to $2.0$ as
energy depletes) produces mean FCR $= 3.099$, which is statistically
significantly higher than the fixed $\alpha = 2.0$ ($p < 0.05$).
This counterintuitive result is explained by the Phase~1 auction
structure: in the fixed $\alpha = 2.0$ scheme, the \emph{anchor}
assignment already uses the quadratic penalty to select a compact, high-probability
first site. The adaptive scheme uses $\alpha = 1.0$ for the Phase~1 anchor
(because the battery is full at sortie start), so the anchor is chosen
with an insufficiently penalized distance, allowing robots to commit to
farther-from-base anchors than the foreclosure argument warrants.
The chain extension then tightens the penalty as energy depletes, but the
sub-optimal anchor placement has already been made.
In other words, the foreclosure cost is front-loaded, not deferred:
the anchor selection is the most consequential decision, and it should
apply the full quadratic penalty even at full battery.

\emph{The fixed $\alpha = 2$ is the theoretically motivated and empirically
robust choice.}
The cost-foreclosure argument of Section~\ref{sec:foreclosure} derives
the quadratic penalty from the superlinear effective cost of anchor
distance at the moment of anchor selection.
The exponent sweep confirms that $\alpha = 2$ sits at the theoretically
predicted optimum, that the loss surface is flat to the right of it (higher
$\alpha$ produces slightly lower point estimates that are not distinguishable
from noise), and that the adaptive bid---which applies a \emph{lower}
exponent than $\alpha = 2$ during the anchor phase---is demonstrably
inferior. The recommendation to use $p/d^2$ therefore rests on both the
theoretical derivation and the empirical exponent sweep.
